%%%%%%%%%%%%%%%%%%%%%%%%%%%%%%%%%%%%%%%%%%%%%%%%%%%%%%%%%%%%%%%%%%%%%%%%%%%%%%
%  PRÉSENTATION BEAMER 16:9 -- gabarit INSA CENTRE VAL DE LOIRE : main.tex
%
%  Ce fichier est réutilisable pour n'importe quelle présentation : il contient le
%  thème (couleurs INSA, page de titre, pied de page, encadrés, styles TikZ,
%  pictogrammes). Le contenu est dans des fichiers séparés, importés avec \include.
%
%  - Nouvelle présentation : modifier le bloc « Informations du document », remplacer
%    la liste des \include (un fichier chap_<nom>.tex par partie du plan ; chaque
%    fichier commence par \section{...} et ne contient ni préambule ni
%    \begin{document}), puis mettre à jour sources.bib.
%  - Une partie du plan = un \section = un fichier. L'ordre des \include donne l'ordre
%    du sommaire. Pour ne compiler que certaines parties (la numérotation est
%    conservée) : décommenter \includeonly{...} plus bas.
%  - Dans une diapositive :
%       \begin{frame}{Titre} ... \end{frame}
%       \intervenant{NOM}   (facultatif) « Intervenant : NOM » en haut à droite
%       \etiquette{Texte}   (facultatif) remplace l'étiquette rouge en haut à gauche,
%                           qui reprend sinon le titre de la section en cours
%       \srcnote{...}       note de source en bas à gauche, ex. \srcnote{Source : \cite{cle}}
%  - Encadrés : idee (rose, bord rouge), pointcle (pêche) et sa variante compacte pcompact
%    (diapositives très remplies), remarque[titre] (gris), encadre{titre} (rouge), carte (beige ;
%    carte[height=3cm,valign=top] pour aligner plusieurs cartes) avec \gros{chiffre ou mot-clé}
%    en grand. Listes compactes dans les encadrés : puces et numeros. Coche / croix en début
%    de ligne : \oui, \non. Tableaux : tabular ou tabularx avec \rowcolor{insarose}.
%  - Ponctuation française : babel-french est désactivé pour « : ; ! ? » (conflit avec TikZ) ;
%    écrire donc texte~: suite, texte~?, texte\,!, texte\,; (espace insécable à la main).
%  - Dans les schémas TikZ d'une diapositive : toujours des accolades autour du corps d'un
%    \foreach (le corps d'une diapositive est lu avant d'être exécuté).
%  - Couleurs : insarouge, insagris (charte INSA), insaorange, insapeche, insapeche2,
%    insabeige, insasafran, insarose (relevées sur la présentation Trove).
%  - Styles TikZ des schémas : pcli, pfoc, psvc, pdat, pgris (pastilles), etiqf,
%    fluxo, fluxr, fluxd, callout, num, acteur, vie, msg, rep (+ \seqmsg, \seqrep).
%  - Pictogrammes TikZ : \pic at (x,y) {icone tel}; liste en fin de préambule.
%  - Page de titre : \insatitleframe (logo images/logo_insa_blanc.png ; à défaut, un
%    texte « INSA » est composé).
%  - Mise en page prévue pour aspectratio=169. Texte de 10 pt : les plus petits textes
%    des schémas restent à \footnotesize (8 pt) ou plus.
%  - Sources : seul le nom de domaine est imprimé dans la liste (package xstring, voir « Bibliographie »),
%    le lien complet reste cliquable dans le PDF. La liste est répartie sur trois diapositives d'après
%    le numéro de citation (bornes dans les filtres srcA, srcB, srcC de la section « Bibliographie »).
%  - Compilation : pdflatex main, biber main, pdflatex main, pdflatex main.
%    (deux passes de pdflatex au minimum : sommaire, décor TikZ de chaque page)
%  - Fichiers du projet : main.tex (gabarit et ordre des parties), un fichier par partie
%    (pour l'instant chap_intro ; les autres parties du plan sont listées, commentées, plus bas),
%    sources.bib, guillotime.sty et images/.
%  - Champs à compléter : chercher \aremplir{...} (surlignage jaune) dans les fichiers de contenu.
%%%%%%%%%%%%%%%%%%%%%%%%%%%%%%%%%%%%%%%%%%%%%%%%%%%%%%%%%%%%%%%%%%%%%%%%%%%%%%
\documentclass[aspectratio=169,10pt,t]{beamer}

%%%%% Informations du document (à modifier pour chaque présentation) %%%%%
\newcommand{\nomproduit}{NIDS}   % nom de travail du projet
\newcommand{\doctitre}{\nomproduit}
\newcommand{\docsoustitre}{Détection d'intrusion réseau\\ par anomalies}
\newcommand{\doctagline}{Apprendre le normal,\\ repérer l'inconnu}
\newcommand{\docauteur}{Groupe jsp le numéro sah}
\newcommand{\docprof}{STI 5A ASL-IA}
\newcommand{\docecole}{INSA Centre Val de Loire}
%\newcommand{\docdate}{\today}
\newcommand{\docdate}{\aujourdhui}

%%%%% Langue, encodage, polices %%%%%
\usepackage[T1]{fontenc}
\usepackage[utf8]{inputenc}
\usepackage{lmodern}          % police machine à écrire (Latin Modern)
\usepackage{tgpagella}        % texte courant : serif proche de celui de la présentation Trove
\usepackage{tgheros}          % titres : sans serif (équivalent d'Arial / Liberation Sans)
% Pas de ligature en police machine à écrire : \texttt{--option} garde ses deux tirets.
\usepackage[protrusion=false,expansion=false,tracking=false,kerning=false,spacing=false]{microtype}
\DisableLigatures[-]{encoding = *, family = tt*}
\usepackage[french]{babel}
% Beamer lit le corps d'une diapositive avant de l'exécuter : les caractères actifs « : ; ! ? » de babel-french
% casseraient alors les coordonnées polaires (\a:1.7) et la syntaxe calc (\p!0.5!\q) de TikZ. On les désactive ;
% les espaces insécables se mettent à la main (texte~: suite, texte~?, texte\,!).
\AtBeginDocument{\shorthandoff{:;!?}}
\usefonttheme{serif}
\usepackage{amsmath}
\usepackage{amssymb}
\usepackage{pifont}           % \ding{51} (coche), \ding{55} (croix)
% symboles courants saisis en UTF-8 dans le texte
\DeclareUnicodeCharacter{2248}{\ensuremath{\approx}}
\DeclareUnicodeCharacter{2192}{\ensuremath{\rightarrow}}
\DeclareUnicodeCharacter{2264}{\ensuremath{\leq}}
\DeclareUnicodeCharacter{2265}{\ensuremath{\geq}}
\DeclareUnicodeCharacter{2212}{\ensuremath{-}}
\DeclareUnicodeCharacter{2260}{\ensuremath{\neq}}

%%%%% Charte graphique INSA CVL %%%%%
\usepackage{xcolor}
\definecolor{insarouge}{HTML}{E42618} % rouge du logo INSA (couleur principale)
\definecolor{insagris}{HTML}{5F5F5F}  % gris du logo INSA (couleur secondaire)
\definecolor{insaorange}{HTML}{C57F17}  % contour et flèches des services « voisins »
\definecolor{insapeche}{HTML}{FCDDB0}   % fond des services « voisins »
\definecolor{insapeche2}{HTML}{FCE2BB}  % fond des encadrés d'idée
% Couleurs relevées sur la page de titre et les encadrés de la présentation Trove
\definecolor{insabeige}{HTML}{F8F0EC}   % panneaux clairs de la page de titre
\definecolor{insasafran}{HTML}{F69F1D}  % blocs orange de la page de titre
\definecolor{insarose}{HTML}{F9CECC}    % fond des encadrés « idée » et des en-têtes de tableau

%%%%% Thème Beamer : structure, marges, listes %%%%%
\setbeamertemplate{navigation symbols}{}
\setbeamersize{text margin left=1.15cm,text margin right=1.15cm}
\setbeamercolor{structure}{fg=insarouge}
\setbeamercolor{normal text}{fg=black,bg=white}
\setbeamercolor{frametitle}{fg=black}
\setbeamercolor{block title}{fg=white,bg=insarouge}
\setbeamercolor{block body}{fg=black,bg=insarose}
\setbeamerfont{frametitle}{family=\sffamily,series=\bfseries,size=\fontsize{15}{18}\selectfont}
\setbeamerfont{framesubtitle}{family=\sffamily,series=\bfseries,size=\small}
\setbeamertemplate{itemize item}{\raisebox{0.05ex}{\textcolor{insarouge}{\footnotesize$\bullet$}}}
\setbeamertemplate{itemize subitem}{\textcolor{insagris}{\textendash}}
\setbeamertemplate{itemize subsubitem}{\textcolor{insagris}{\tiny$\bullet$}}
\setbeamertemplate{enumerate item}{\textcolor{insarouge}{\bfseries\insertenumlabel.}}
\setbeamerfont{itemize/enumerate body}{size=\normalsize}
\setbeamerfont{itemize/enumerate subbody}{size=\small}
\setlength{\leftmargini}{1.1em}
\setlength{\leftmarginii}{1.1em}
\addtolength{\itemsep}{0.1em}
\hypersetup{pdftitle={\doctitre}, pdfauthor={\docauteur}}

%%%%% Décor de chaque diapositive (étiquette de section, intervenant, pied de page rouge) %%%%%
% Unités du décor : 1/935 de la largeur et 1/527 de la hauteur de la page (mesures relevées sur la
% présentation Trove). \etiquette{...} et \intervenant{...} s'utilisent à l'intérieur d'une diapositive.
\usepackage{etoolbox}
\def\insatag{\insertsection}
\def\insainterv{}
\newcommand{\etiquette}[1]{\gdef\insatag{#1}}
\newcommand{\intervenant}[1]{\gdef\insainterv{#1}}
% les deux réglages ne valent que pour la diapositive en cours : remis à zéro au début de chaque frame
\AtBeginEnvironment{frame}{\gdef\insatag{\insertsection}\gdef\insainterv{}}
\newcommand{\insaunitx}{0.0010695\paperwidth}
\newcommand{\insaunity}{0.0018975\paperheight}
\setbeamertemplate{background}{%
  \begin{tikzpicture}[remember picture,overlay,shift={(current page.north west)}]
    \node[anchor=north west,inner sep=0pt,text=insarouge,font=\sffamily\bfseries\scriptsize]
      at (1.15cm,-0.28cm) {\insatag};
    \ifx\insainterv\empty\else
      \node[anchor=north east,inner sep=0pt,text=insagris!70,font=\scriptsize]
        at (\paperwidth-1.15cm,-0.28cm) {Intervenant : \insainterv};
    \fi
    \begin{scope}[shift={(current page.south east)},x=\insaunitx,y=-\insaunity]
      \foreach \i in {0,...,17} {%
        \draw[insarouge,line width={0.3pt+0.3pt*mod(\i,2)}]
          ({-55+3.0*\i+1},-109) -- ({-55+3.0*\i+1},-55);}
      \foreach \i in {0,...,4} {%
        \draw[insagris!55,line width=0.35pt] ({-122+11*\i},-55) -- ({-122+11*\i},0);}
      \fill[insarouge] (-55,-55) rectangle (0,0);
      \node[white,font=\bfseries\large] at (-27.5,-27.5) {\insertframenumber};
    \end{scope}
  \end{tikzpicture}}
\setbeamertemplate{headline}{}
% marge basse de 0,62 cm réservée aux notes de source (\srcnote) : le contenu ne descend pas plus bas
\setbeamertemplate{footline}{\hbox{\vrule width 0pt height 0pt depth 0.62cm}}
\setbeamertemplate{frametitle}{%
  \vskip0.52cm
  {\usebeamerfont{frametitle}\usebeamercolor[fg]{frametitle}\insertframetitle\par}%
  \vskip0.08cm}
\setbeamertemplate{framesubtitle}{%
  {\usebeamerfont{framesubtitle}\usebeamercolor[fg]{frametitle}\insertframesubtitle\par}%
  \vskip0.05cm}

% Note de source en bas à gauche de la diapositive : \srcnote{Sources : \cite{cle}}
\newcommand{\srcnote}[1]{%
  \begin{tikzpicture}[remember picture,overlay]
    \node[anchor=south west,inner sep=0pt,text=insagris,font=\tiny,text width=0.72\paperwidth,align=left]
      at ([xshift=1.15cm,yshift=0.22cm]current page.south west) {#1};
  \end{tikzpicture}}

%%%%% Sommaire : numéros rouges, pas de sous-parties %%%%%
\setbeamertemplate{section in toc}{%
  \leavevmode\strut\textcolor{insarouge}{\bfseries\inserttocsectionnumber.}\hspace{0.6em}\inserttocsection\strut\par\vspace{-0.85em}}
\setbeamerfont{section in toc}{size=\large}
\setbeamercolor{section in toc}{fg=black}

%%%%% Tableaux, graphiques, schémas %%%%%
\usepackage{array}
\usepackage{makecell}
\usepackage{booktabs}
\usepackage{colortbl}
\usepackage{tabularx}
\usepackage{multirow}
\usepackage{graphicx}
\graphicspath{{images/}}
\usepackage{tikz}
\usetikzlibrary{positioning, arrows.meta, fit, backgrounds, calc, shapes.geometric, babel,
                decorations.pathmorphing, decorations.pathreplacing, patterns}
\usepackage{pgfplots}
\pgfplotsset{compat=1.16}
\usepackage[skins]{tcolorbox}
\tcbset{boxrule=0.6pt,arc=2pt,left=6pt,right=6pt,top=4pt,bottom=4pt}

%%%%% Pour la date %%%%%
\usepackage[delan,parigot]{guillotime}

%%%%% Lignes de code (frames [fragile]) %%%%%
\usepackage{listings}
\lstset{
  basicstyle=\ttfamily\footnotesize,
  backgroundcolor=\color{insagris!8},
  frame=single,
  rulecolor=\color{insagris!60},
  breaklines=true,
  upquote=true,
  showstringspaces=false,
  literate={é}{{\'e}}1 {è}{{\`e}}1 {ê}{{\^e}}1 {ë}{{\"e}}1 {à}{{\`a}}1 {â}{{\^a}}1
           {î}{{\^i}}1 {ï}{{\"i}}1 {ô}{{\^o}}1 {ù}{{\`u}}1 {û}{{\^u}}1 {ç}{{\c c}}1
           {É}{{\'E}}1 {È}{{\`E}}1 {À}{{\`A}}1 {Ç}{{\c C}}1 {œ}{{\oe}}1
}

%%%%% Encadrés %%%%%
% Idée à retenir (rose, bord rouge) ; point clé en une phrase (pêche, sans bord) ; remarque grise ; encadré titré.
% (« definition » existe déjà dans Beamer : l'encadré titré s'appelle donc « encadre ».)
\newtcolorbox{idee}{colback=insarose,colframe=insarouge,boxrule=0.7pt,arc=2pt}
\newtcolorbox{pointcle}{colback=insapeche2,colframe=insapeche2,boxrule=0pt,arc=3pt}
% Variante compacte du point clé (marges haute et basse réduites), pour les diapositives très remplies.
\newtcolorbox{pcompact}{colback=insapeche2,colframe=insapeche2,boxrule=0pt,arc=3pt,top=2pt,bottom=2pt}
\newtcolorbox{remarque}[1][Remarque]{colback=insagris!7,colframe=insagris,
  fonttitle=\bfseries\sffamily\small,title={#1}}
\newtcolorbox{encadre}[1]{colback=insarouge!4,colframe=insarouge,
  fonttitle=\bfseries\sffamily\small,title={#1},before upper={\raggedright}}
% Carte beige (chiffre clé, vignette) ; \gros{...} : chiffre ou mot-clé en grand, rouge
\newtcolorbox{carte}[1][]{colback=insabeige,colframe=insabeige,boxrule=0pt,arc=3pt,
  left=5pt,right=5pt,top=3pt,bottom=3pt,before upper={\raggedright},#1}
\newcommand{\gros}[1]{{\sffamily\bfseries\large\color{insarouge}#1}}
% \enorme{...} : chiffre d'accroche, très grand (diapositives d'ouverture)
\newcommand{\enorme}[1]{{\sffamily\bfseries\fontsize{26}{28}\selectfont\color{insarouge}#1}}
% \aremplir{...} : information manquante, surlignée en jaune pour ne pas passer inaperçue
\newcommand{\aremplir}[1]{\colorbox{yellow!60}{\strut\textbf{[#1]}}}
% Coche rouge / croix grise en début de ligne : \oui Texte ...   \non Texte ...
\newcommand{\oui}{\textcolor{insarouge}{\ding{51}}\hspace{0.35em}}
\newcommand{\non}{\textcolor{insagris}{\ding{55}}\hspace{0.35em}}
% Listes compactes (corps en \footnotesize) pour les encadrés : \begin{puces}\item ...\end{puces}
% et \begin{numeros}\item ...\end{numeros}  (les listes beamer normales restent en taille normale)
\newenvironment{puces}
  {\begingroup\setbeamerfont{itemize/enumerate body}{size=\footnotesize}\begin{itemize}\setlength{\itemsep}{1pt}}
  {\end{itemize}\endgroup}
\newenvironment{numeros}
  {\begingroup\setbeamerfont{itemize/enumerate body}{size=\footnotesize}\begin{enumerate}\setlength{\itemsep}{1pt}}
  {\end{enumerate}\endgroup}

%%%%% Styles TikZ pour les schémas (repris du gabarit de cours) %%%%%
% svc    : boîte d'un composant rouge pâle          svcgris : composant secondaire
% bus    : barre horizontale                         flux    : flèche entre deux composants
\tikzset{
  svc/.style     = {draw=insarouge, fill=insarouge!8, thick, rounded corners=2pt,
                    minimum height=7mm, inner sep=3pt, align=center},
  svcgris/.style = {draw=insagris, fill=insagris!8, thick, rounded corners=2pt,
                    minimum height=7mm, inner sep=3pt, align=center},
  bus/.style     = {draw=insagris, fill=insagris!15, thick, minimum height=6mm, align=center},
  flux/.style    = {-{Stealth[length=2mm]}, thick, insagris},
}
% Schémas « à la Trove » : pastilles colorées par rôle, flèches étiquetées, encadré d'idée.
%   pcli  : client / utilisateur                  pfoc : élément au centre du schéma (rouge plein)
%   psvc  : élément voisin (pêche)                pdat : donnée, ressource (rose)
%   pgris : infrastructure, existant              etiqf: étiquette posée sur une flèche
%   fluxo / fluxr / fluxd : flèche orange / rouge / rouge pointillée
%   callout : encadré pêche sous le schéma         num : pastille numérotée d'une étape
% Diagrammes de séquence : acteur (en-tête), vie (ligne de vie), msg / rep (message / réponse)
\tikzset{
  pastille/.style = {rounded corners=0.3cm, minimum height=0.65cm, inner xsep=7pt, inner ysep=3pt,
                     thick, align=center, font=\footnotesize\bfseries},
  pcli/.style  = {pastille, draw=insarouge, fill=white, text=insarouge},
  pfoc/.style  = {pastille, draw=black, fill=insarouge, text=white, line width=1.2pt},
  psvc/.style  = {pastille, draw=insaorange, fill=insapeche, text=black},
  pdat/.style  = {pastille, draw=insarouge!90!black, fill=insarouge!30, text=black},
  pgris/.style = {pastille, draw=insagris, fill=insagris!12, text=black},
  etiqf/.style = {fill=white, inner sep=1.5pt, font=\scriptsize, text=insagris, align=center},
  fluxo/.style = {-{Stealth[length=2mm]}, thick, insaorange},
  fluxr/.style = {-{Stealth[length=2mm]}, thick, insarouge},
  fluxd/.style = {-{Stealth[length=2mm]}, thick, insarouge, dashed},
  callout/.style = {fill=insapeche2, rounded corners=3pt, inner sep=5pt, align=center,
                    font=\footnotesize},
  num/.style   = {circle, fill=insarouge, text=white, font=\scriptsize\bfseries, inner sep=1pt,
                  minimum size=4.5mm},
  acteur/.style = {draw=insarouge, fill=insarouge!8, thick, rounded corners=2pt, minimum height=0.7cm,
                   minimum width=2.2cm, align=center, font=\footnotesize\bfseries, inner sep=3pt},
  vie/.style   = {insagris!70, dashed, thin},
  msg/.style   = {-{Stealth[length=2mm]}, thick, insagris},
  rep/.style   = {-{Stealth[length=2mm]}, thick, insagris, dashed},
}
% Messages d'un diagramme de séquence : \seqmsg{x de départ}{x d'arrivée}{y}{numéro (ou vide)}{étiquette}
% (flèche pleine) et \seqrep (flèche pointillée = réponse).
\newcommand{\seqmsg}[5]{%
  \draw[msg] (#1,#3) -- (#2,#3) node[midway, above=1pt, font=\footnotesize, align=center, text=black, fill=white, inner sep=1pt] {#5};%
  \ifx\relax#4\relax\else\node[num] at ({#1+(#2>#1?-0.4:0.4)},#3) {#4};\fi}
\newcommand{\seqrep}[5]{%
  \draw[rep] (#1,#3) -- (#2,#3) node[midway, above=1pt, font=\footnotesize, align=center, text=black, fill=white, inner sep=1pt] {#5};%
  \ifx\relax#4\relax\else\node[num] at ({#1+(#2>#1?-0.4:0.4)},#3) {#4};\fi}

%%%%% Pictogrammes TikZ : \pic at (x,y) {icone ...}; (centrés, environ 1 cm de côté) %%%%%
% Ils prennent la couleur courante : \pic[insarouge] at (0,0) {icone tel}; \pic[scale=1.5] ...
% Liste : tel, tour, antenne, mont, immeuble, vague, cadenas, euro, code, usine, cairn,
%         croix, coche, tente, boite, personne, drapeau, batterie, pin, alerte
\tikzset{
  pics/icone tel/.style = {code={
    \draw[thick,fill=white,rounded corners=1.5pt] (-0.2,-0.36) rectangle (0.2,0.36);
    \draw[thin,fill=insagris!12] (-0.145,-0.26) rectangle (0.145,0.27);
    \fill (0,-0.31) circle (0.025);}},
  pics/icone tour/.style = {code={
    \draw[thick] (0,0.3) -- (-0.18,-0.38) (0,0.3) -- (0.18,-0.38)
                 (-0.1,-0.1) -- (0.1,-0.1) (-0.14,-0.24) -- (0.14,-0.24);
    \fill (0,0.32) circle (0.04);
    \foreach \r in {0.17,0.3} {
      \draw[thin] (0,0.32) ++(-35:\r) arc (-35:35:\r);
      \draw[thin] (0,0.32) ++(145:\r) arc (145:215:\r);}}},
  pics/icone antenne/.style = {code={
    \draw[thick,fill=white,rounded corners=1.5pt] (-0.34,-0.36) rectangle (0.3,0.0);
    \fill[insagris!35] (-0.27,-0.3) rectangle (-0.02,-0.08);
    \fill (0.12,-0.1) circle (0.035);
    \draw[thick] (0.2,0.0) -- (0.2,0.34);
    \fill (0.2,0.37) circle (0.035);
    \foreach \r in {0.14,0.24} {
      \draw[thin] (0.2,0.37) ++(-40:\r) arc (-40:40:\r);
      \draw[thin] (0.2,0.37) ++(140:\r) arc (140:220:\r);}}},
  pics/icone mont/.style = {code={
    \fill[insagris!40] (-0.5,-0.3) -- (-0.15,0.28) -- (0.2,-0.3) -- cycle;
    \fill[insagris] (-0.12,-0.3) -- (0.24,0.4) -- (0.6,-0.3) -- cycle;
    \fill[white] (0.24,0.4) -- (0.13,0.18) -- (0.2,0.2) -- (0.25,0.14) -- (0.31,0.2) -- (0.37,0.18) -- cycle;}},
  pics/icone immeuble/.style = {code={
    \draw[thick,fill=white] (-0.34,-0.36) -- (-0.34,0.22) -- (-0.18,0.38) -- (-0.08,0.16) -- (0.06,0.3)
      -- (0.14,0.06) -- (0.3,0.2) -- (0.3,-0.36) -- cycle;
    \foreach \x in {-0.2,0.02,0.18} \foreach \y in {-0.08,-0.25}
      \fill (\x-0.04,\y-0.04) rectangle (\x+0.04,\y+0.04);}},
  pics/icone vague/.style = {code={
    \foreach \y in {-0.22,0,0.22}
      \draw[thick,domain=-0.4:0.4,samples=24,smooth] plot (\x,{\y+0.06*sin(deg(\x*15))});}},
  pics/icone cadenas/.style = {code={
    \draw[thick] (-0.12,0.02) -- (-0.12,0.14) arc (180:0:0.12) -- (0.12,0.02);
    \draw[thick,fill=white,rounded corners=1pt] (-0.2,-0.32) rectangle (0.2,0.02);
    \fill (0,-0.15) circle (0.035);}},
  pics/icone euro/.style = {code={
    \draw[thick,fill=white] (0,0) circle (0.33);
    \node[font=\bfseries\small,text=black] at (0,0) {\texteuro};}},
  pics/icone code/.style = {code={
    \draw[thick,line cap=round,line join=round] (-0.1,0.2) -- (-0.32,0) -- (-0.1,-0.2)
         (0.1,0.2) -- (0.32,0) -- (0.1,-0.2) (0.07,0.27) -- (-0.07,-0.27);}},
  pics/icone usine/.style = {code={
    \draw[thick,fill=white] (-0.4,-0.3) -- (-0.4,0.06) -- (-0.2,0.2) -- (-0.2,0.06) -- (0,0.2) -- (0,0.06)
      -- (0.2,0.2) -- (0.2,-0.3) -- cycle;
    \draw[thick,fill=white] (0.2,-0.3) rectangle (0.34,0.36);
    \fill (-0.3,-0.12) rectangle (-0.22,-0.04);\fill (-0.1,-0.12) rectangle (-0.02,-0.04);
    \fill (0.07,-0.12) rectangle (0.15,-0.04);}},
  pics/icone cairn/.style = {code={
    \fill[insagris!70] (0,-0.34) ellipse (0.46 and 0.14);
    \fill[insagris!45] (0.03,-0.1) ellipse (0.36 and 0.12);
    \fill[insarouge]   (-0.02,0.1) ellipse (0.27 and 0.11);
    \fill[insagris!30] (0.02,0.27) ellipse (0.18 and 0.085);
    \fill[insasafran]  (0,0.4) ellipse (0.1 and 0.065);}},
  pics/icone croix/.style = {code={
    \fill[insarouge] (0,0) circle (0.17);
    \draw[white,line width=1.3pt,line cap=round] (-0.07,-0.07) -- (0.07,0.07) (-0.07,0.07) -- (0.07,-0.07);}},
  pics/icone coche/.style = {code={
    \fill (0,0) circle (0.17);
    \draw[white,line width=1.3pt,line cap=round,line join=round] (-0.08,0.0) -- (-0.025,-0.06) -- (0.08,0.07);}},
  pics/icone tente/.style = {code={
    \draw[thick,fill=white] (-0.42,-0.28) -- (0,0.32) -- (0.42,-0.28) -- cycle;
    \fill (-0.11,-0.28) -- (0,0.0) -- (0.11,-0.28) -- cycle;
    \draw[thick] (-0.5,-0.28) -- (0.5,-0.28);}},
  pics/icone boite/.style = {code={
    \draw[thick,fill=white] (-0.3,-0.3) rectangle (0.2,0.1);
    \draw[thick,fill=insagris!15] (-0.3,0.1) -- (-0.15,0.3) -- (0.35,0.3) -- (0.2,0.1) -- cycle;
    \draw[thick,fill=insagris!30] (0.2,0.1) -- (0.35,0.3) -- (0.35,-0.1) -- (0.2,-0.3) -- cycle;}},
  pics/icone personne/.style = {code={
    \fill (0,0.22) circle (0.1);
    \fill (-0.2,-0.32) -- (-0.2,-0.14) arc (180:0:0.2) -- (0.2,-0.32) -- cycle;}},
  pics/icone drapeau/.style = {code={
    \fill[blue!60!black] (-0.3,-0.2) rectangle (-0.1,0.2);
    \fill[white] (-0.1,-0.2) rectangle (0.1,0.2);
    \fill[insarouge] (0.1,-0.2) rectangle (0.3,0.2);
    \draw[insagris!60,thin] (-0.3,-0.2) rectangle (0.3,0.2);}},
  pics/icone batterie/.style = {code={
    \draw[thick,fill=white,rounded corners=1pt] (-0.34,-0.18) rectangle (0.26,0.18);
    \fill (0.26,-0.07) rectangle (0.33,0.07);
    \fill (-0.27,-0.11) rectangle (-0.07,0.11);\fill (-0.03,-0.11) rectangle (0.17,0.11);}},
  pics/icone pin/.style = {code={
    \fill (0,0.12) circle (0.2);
    \fill (-0.19,0.06) -- (0,-0.36) -- (0.19,0.06) -- cycle;
    \fill[white] (0,0.12) circle (0.08);}},
  pics/icone alerte/.style = {code={
    \draw[thick,fill=insapeche,line join=round] (0,0.34) -- (-0.4,-0.3) -- (0.4,-0.3) -- cycle;
    \node[font=\bfseries\small,text=black,inner sep=0pt] at (0,-0.08) {!};}},
}

%%%%% Page de titre (décor de la page de titre de la présentation Trove) %%%%%
% \insatitleframe : diapositive 1, sans pied de page ; logo, titre, sous-titre, auteur, école, date.
\newcommand{\insalogo}[1]{%
  \IfFileExists{images/logo_insa_blanc.png}%
    {\includegraphics[width=#1]{images/logo_insa_blanc.png}}%
    {\textcolor{white}{\sffamily\bfseries\Large INSA}\ \textcolor{white}{\sffamily\tiny CENTRE VAL DE LOIRE}}}
\newcommand{\insatitleframe}{%
  \begingroup
  \setbeamertemplate{background}{}%
  \begin{frame}[plain]
    \begin{tikzpicture}[remember picture,overlay,shift={(current page.north west)},x=\insaunitx,y=-\insaunity]
      % panneaux beiges, bloc rouge, blocs orange
      \fill[insabeige] (669,0) rectangle (935,295);
      \fill[insabeige] (609,235) rectangle (669,295);
      \fill[insabeige] (728,295) rectangle (935,527);
      \fill[insarouge] (0,0) -- (669,0) -- (669,235) -- (609,235) -- (609,295) -- (550,295) -- (550,353)
        -- (0,353) -- cycle;
      \fill[insasafran] (609,295) rectangle (728,353);
      \fill[insasafran] (609,353) rectangle (669,413);
      % lignes décoratives
      \foreach \i in {0,...,30} {\draw[insarouge,line width=0.3pt] ({842+3.0*\i},0) -- ({842+3.0*\i},295);}
      \foreach \x in {749,772,796,819} {\draw[insasafran,line width=0.4pt] (\x,295) -- (\x,527);}
      \draw[insagris!60,line width=0.4pt] (841,295) -- (841,527);
      \foreach \x in {65,76,87,98,109} {\draw[insarouge,line width=0.5pt] (\x,476) -- (\x,527);}
      % grille de « + »
      \foreach \i in {0,...,7} \foreach \j in {0,...,7} {%
        \draw[white,line width=0.4pt] ({503+21.1*\i-3},{16+21.1*\j}) -- ({503+21.1*\i+3},{16+21.1*\j})
                                      ({503+21.1*\i},{16+21.1*\j-3}) -- ({503+21.1*\i},{16+21.1*\j+3});}
      % logo, titre, sous-titre
      \node[anchor=north west,inner sep=0pt] at (65,58) {\insalogo{0.1647\paperwidth}};
      \node[anchor=north west,inner sep=0pt,text=white,text width=0.55\paperwidth,align=left] at (63,170)
        {{\sffamily\bfseries\fontsize{32}{34}\selectfont\doctitre\par}\vspace{0.25cm}%
         {\sffamily\fontsize{14}{17}\selectfont\docsoustitre\par}};
      % auteur, cursus, école, date
      \node[anchor=north west,inner sep=0pt,align=left,text width=0.62\paperwidth] at (63,392)
        {{\sffamily\fontsize{12}{15}\selectfont\textbf{\docauteur}\ ---\ \docprof\par}%
         {\sffamily\fontsize{10}{14}\selectfont\color{insagris}\docecole\ ---\ \docdate\par}};
      \node[font=\bfseries\large] at (908,497) {\insertframenumber};
    \end{tikzpicture}
  \end{frame}
  \endgroup}

%%%%% Bibliographie (sources.bib) %%%%%
\usepackage{csquotes}
\usepackage[backend=biber,style=numeric-comp,sorting=none,maxbibnames=3,minbibnames=1,
            doi=false,isbn=false,eprint=false]{biblatex}
\addbibresource{sources.bib}
\setcounter{biburlnumpenalty}{100}
\setcounter{biburlucpenalty}{100}
\setcounter{biburllcpenalty}{100}
\renewcommand*{\bibfont}{\fontsize{6}{7.2}\selectfont\raggedright}
\setlength{\bibitemsep}{0.1\baselineskip}
\AtEveryBibitem{\clearfield{urlyear}\clearfield{urlmonth}\clearfield{urlday}}
\urlstyle{same}
% Adresse imprimée = nom de domaine seulement (xstring) ; le lien complet reste cliquable dans le PDF.
% Cela permet de faire tenir de nombreuses sources sur une seule diapositive.
\usepackage{xstring}
\newcommand{\bibhost}[1]{\StrBehind{#1}{://}[\bibhosttmp]\StrBefore{\bibhosttmp}{/}[\bibhosttmpb]%
  \IfBeginWith{\bibhosttmpb}{www.}{\StrGobbleLeft{\bibhosttmpb}{4}}{\bibhosttmpb}}
\DeclareFieldFormat{url}{\href{#1}{\bibhost{#1}}}
\usepackage{multicol}
% La liste des sources est coupée en trois diapositives d'après le numéro de citation (ordre de première citation).
% \srcrange{n1}{n2} garde les sources numérotées de n1 à n2 ; ajuster les bornes si le nombre de sources change.
\newcommand{\srcrange}[2]{%
  \iffieldint{labelnumber}%
    {\ifboolexpe{test{\ifnumless{\thefield{labelnumber}}{#1}} or test{\ifnumgreater{\thefield{labelnumber}}{#2}}}{\skipentry}{}}%
    {\skipentry}}
\defbibcheck{srcA}{\srcrange{1}{25}}
\defbibcheck{srcB}{\srcrange{26}{50}}
\defbibcheck{srcC}{\srcrange{51}{999}}

\title{\doctitre}
\author{\docauteur}
\date{\docdate}

% Pour ne compiler que certaines parties (la numérotation est conservée) :
%\includeonly{chap_finances}

\begin{document}

%%%%% Page de titre : diapositive 1 %%%%%
\insatitleframe

%%%%% Sommaire (deux colonnes : parties 1 à 4, puis 5 à 8 ; ajuster si le nombre de parties change) %%%%%
% Le sommaire se remplit au fur et à mesure que les \include ci-dessous sont activés.
\begin{frame}{Sommaire}
  \etiquette{Plan}
  \begin{columns}[T]
    \begin{column}{0.37\textwidth}
      \vspace{0.2cm}
      \tableofcontents[sections={1-4},hideallsubsections]
    \end{column}
    \begin{column}{0.37\textwidth}
      \vspace{0.2cm}
      \tableofcontents[sections={5-8},hideallsubsections]
    \end{column}
    \begin{column}{0.22\textwidth}
      \centering
      \vspace{0.3cm}
      \begin{tikzpicture}[scale=2.3]
        \pic at (0,0) {icone alerte};
      \end{tikzpicture}\\[0.2cm]
      {\sffamily\bfseries\color{insarouge}\doctagline}
    \end{column}
  \end{columns}
\end{frame}

%%%%% Parties du plan (un fichier par partie, chacun commence par \section) %%%%%
% L'ordre des \include donne l'ordre du sommaire : pour déplacer une partie, déplacer sa ligne.
\section{Introduction}

\begin{frame}{Les attaquants vont plus vite que les correctifs}
  \begin{columns}[T]
    \begin{column}{0.48\textwidth}
      \begin{carte}[height=3.6cm,valign=top]
        \enorme{7 jours}\\[0.25cm]
        Délai moyen d'exploitation d'une vulnérabilité en 2025 (estimation)~: l'attaque
        précède souvent la publication du correctif.
      \end{carte}
    \end{column}
    \begin{column}{0.48\textwidth}
      \begin{carte}[height=3.6cm,valign=top]
        \enorme{22 secondes}\\[0.25cm]
        Délai médian entre l'accès initial et sa passation à un autre groupe criminel
        en 2025, contre plus de 8~heures en 2022.
      \end{carte}
    \end{column}
  \end{columns}
  \vspace{0.45cm}
  \begin{pointcle}
    \centering\large
    Un détecteur qui ne reconnaît que ce qu'il a déjà vu arrive trop tard.
  \end{pointcle}
  \srcnote{Source~: \cite{mtrends26}}
\end{frame}

\begin{frame}{Et pourtant, la détection reste lente}
  \begin{columns}[T]
    \begin{column}{0.31\textwidth}
      \begin{carte}[height=3.4cm,valign=top]
        \enorme{14}\,\gros{jours}\\[0.25cm]
        Durée de présence médiane avant détection (11~jours en 2024).
        Espionnage~: 122~jours.
      \end{carte}
    \end{column}
    \begin{column}{0.31\textwidth}
      \begin{carte}[height=3.4cm,valign=top]
        \enorme{241}\,\gros{jours}\\[0.25cm]
        Moyenne pour identifier et contenir une violation (2025).
        Mesure différente de la médiane de Mandiant.
      \end{carte}
    \end{column}
    \begin{column}{0.31\textwidth}
      \begin{carte}[height=3.4cm,valign=top]
        \enorme{−80}\,\gros{jours}\\[0.25cm]
        Gagnés en moyenne, et 1,9~M\$ économisés, avec une IA et une automatisation
        très utilisées.
      \end{carte}
    \end{column}
  \end{columns}
  \vspace{0.4cm}
  \begin{pointcle}
    \centering
    Détecter plus tôt, c'est payer moins~: l'IA est associée à des violations plus courtes.
  \end{pointcle}
  \srcnote{Sources~: \cite{mtrends26,ibm25}}
\end{frame}

\begin{frame}{Deux philosophies de détection}
  \begin{columns}[T]
    \begin{column}{0.48\textwidth}
      \begin{encadre}{Par signatures}
        \textbf{Principe~:} reconnaître un motif d'attaque déjà connu.\\[0.25cm]
        \oui Peu de fausses alertes sur les menaces répertoriées\\[0.1cm]
        \non Aveugle aux attaques inconnues (zero-day, variantes)
      \end{encadre}
    \end{column}
    \begin{column}{0.48\textwidth}
      \begin{encadre}{Par anomalies \hfill \textit{notre sujet}}
        \textbf{Principe~:} apprendre le trafic normal, signaler les écarts.\\[0.25cm]
        \oui Peut détecter l'inconnu, même sans données étiquetées\\[0.1cm]
        \non Fausses alertes fréquentes~; le «~normal~» évolue
      \end{encadre}
    \end{column}
  \end{columns}
  \vspace{0.4cm}
  \begin{pointcle}
    \textbf{Exemple~: Kitsune} (NDSS 2018), un ensemble d'autoencodeurs qui apprend en ligne,
    sans supervision, à distinguer trafic normal et anormal, et tourne même sur un Raspberry~Pi.
  \end{pointcle}
  \srcnote{Sources~: \cite{survey07,kitsune18}}
\end{frame}

\begin{frame}{Une promesse difficile à tenir en production}
  Dès 2010, Sommer et Paxson relevaient un écart frappant entre la recherche et les
  déploiements réels. Quatre obstacles reviennent encore.
  \vspace{0.25cm}
  \begin{columns}[T]
    \begin{column}{0.48\textwidth}
      \begin{carte}[height=2.15cm,valign=top]
        \gros{Faux positifs}\\[0.05cm]
        Trop de fausses alertes~: les vraies sont noyées (\emph{alert fatigue}).
      \end{carte}
      \vspace{0.15cm}
      \begin{carte}[height=2.15cm,valign=top]
        \gros{Trafic chiffré}\\[0.05cm]
        Environ 98,6\,\% du trafic web humain est chiffré (2026)~: le contenu des paquets
        est illisible.
      \end{carte}
    \end{column}
    \begin{column}{0.48\textwidth}
      \begin{carte}[height=2.15cm,valign=top]
        \gros{Dérive du trafic}\\[0.05cm]
        Nouvelle application, mise à jour~: le «~normal~» change, le modèle se périme.
      \end{carte}
      \vspace{0.15cm}
      \begin{carte}[height=2.15cm,valign=top]
        \gros{Qualité des données}\\[0.05cm]
        CIC-IDS2017, référence du domaine~: plus de 20\,\% des traces reconstruites ou
        ré-étiquetées.
      \end{carte}
    \end{column}
  \end{columns}
  \srcnote{Sources~: \cite{sommer10,survey07,zanetti15,cloudflare26,engelen21,lanvin22}}
\end{frame}

\begin{frame}{Problématique et périmètre}
  \begin{idee}
    \centering\large
    Comment rendre un NIDS par anomalies fiable malgré des données réseau imparfaites~?
  \end{idee}
  \vspace{0.3cm}
  \begin{columns}[T]
    \begin{column}{0.58\textwidth}
      \begin{carte}[height=3.1cm,valign=top]
        \gros{Notre démarche}\\[0.1cm]
        \begin{itemize}
          \item Supervisé ou non supervisé~?
          \item Un modèle, évalué sur cicids2017
          \item Augmentation de données (SMOTE)
          \item Des démonstrations
        \end{itemize}
      \end{carte}
    \end{column}
    \begin{column}{0.38\textwidth}
      \begin{carte}[height=3.1cm,valign=top]
        \gros{Périmètre}\\[0.1cm]
        \oui Détection par anomalies sur le trafic réseau\\[0.1cm]
        \non Attaques utilisant l'IA\\[0.05cm]
        \non Défense des modèles d'IA eux-mêmes
      \end{carte}
    \end{column}
  \end{columns}
  \vspace{0.3cm}
  \begin{pcompact}
    \centering
    Première étape~: supervisé ou non supervisé, que change le choix pour un NIDS~?
  \end{pcompact}
\end{frame}

\section{Supervisé ou non supervisé}

\begin{frame}{partie 2}
  \vspace{0.45cm}
  \begin{pointcle}
    \centering\large
    Recherche sur les différences entre l'apprentissage supervisé et non suoervisé dans notre contexte. E
  \end{pointcle}
\end{frame}


\section{Le modèle}

\begin{frame}{partie 3}
  \vspace{0.45cm}
  \begin{pointcle}
    \centering\large
    Présentation du modèle. M
  \end{pointcle}
\end{frame}


\section{Le Dataset}

\begin{frame}{partie 4}
  \vspace{0.45cm}
  \begin{pointcle}
    \centering\large
    Présentation du Dataset plus précis (annoncé en intro) → Transition sur augmentation de données. Y
  \end{pointcle}
\end{frame}


\section{Augmentation de Données - SMOTE}

\begin{frame}{partie 5}
  \vspace{0.45cm}
  \begin{pointcle}
    \centering\large
    Présentation de l'Augmentation de données (2+) N + H
  \end{pointcle}
\end{frame}


\section{Démonstrations}

\begin{frame}{partie 6}
  \vspace{0.45cm}
  \begin{pointcle}
    \centering\large
    Démos (2+) M
  \end{pointcle}
\end{frame}


\section{Hyperspécialisation des algorithmes et données}

\begin{frame}{partie 7}
  \vspace{0.45cm}
  \begin{pointcle}
    \centering\large
    Hyperspécilisation des algos et des modèles ? G
  \end{pointcle}
\end{frame}


\section{Partie Intéractive}

\begin{frame}{partie 8}
  \vspace{0.45cm}
  \begin{pointcle}
    \centering\large
    Partie Interractive ??????? M E
  \end{pointcle}
\end{frame}




%%%%% Sources (sources.bib : appel \cite{cle} dans les diapositives) %%%%%
% Liste sur 3 colonnes, un peu plus large que le texte courant, coupée en trois diapositives (filtres srcA, srcB, srcC).
\newcommand{\diaposources}[2]{%
  \begin{frame}{#2}
    \etiquette{Sources}
    \makebox[\textwidth][c]{\begin{minipage}{13.6cm}
      \setlength{\columnsep}{0.3cm}
      \begin{multicols}{3}
        \printbibliography[heading=none,check=#1]
      \end{multicols}
    \end{minipage}}
  \end{frame}}
% Une seule diapositive de sources pour l'instant (moins de 25 références). Quand il y en aura plus,
% remettre le titre « Sources 1/3 » et décommenter les deux lignes suivantes.
\diaposources{srcA}{Sources~: pages consultées au \aujourdhui}
%\diaposources{srcB}{Sources 2/3}
%\diaposources{srcC}{Sources 3/3}

\end{document}
